\medskip

L'espace est rapporté à un repère orthonormé \Oijk.

On considère les points 

\begin{center}A$(4~;~-1~;~3)$,\quad B$(-1~;~1~;~-2)$,\quad C(0~;~4~;~5) et D$(-3~;~-4~;~6)$.\end{center}

\begin{enumerate}
\item
	\begin{enumerate}
		\item Vérifier que les points A, B, C ne sont pas alignés.
	\end{enumerate}

On admet qu’une équation cartésienne du plan (ABC) est : $29x + 30y - 17z = 35$.

	\begin{enumerate}[resume]
		\item Les points A, B, C, D sont-ils coplanaires ? Justifier.
	\end{enumerate}
\end{enumerate}

On admet que lorsque quatre points ne sont pas coplanaires, il existe un unique point situé à égale distance de ces quatre points.

L’objectif de cet exercice est de déterminer le point H se situant à égale distance des quatre points A, B, C, D.

On définit le plan médiateur d’un segment comme le plan passant par le milieu de ce segment et orthogonal à la droite portant ce segment. C’est l’ensemble des points équidistants des extrémités de ce segment.

\begin{enumerate}[resume]
\item Soit $P_1$ le plan médiateur du segment [AB].
	\begin{enumerate}
		\item Déterminer les coordonnées du milieu du segment [AB].
		\item En déduire qu'une équation cartésienne de $P_1$ est : $5x - 2y + 5z = 10$.
	\end{enumerate}
\item On note $P_2$ le plan médiateur du segment [CD].
	\begin{enumerate}
		\item Soit M un point du plan $P_2$ de coordonnées $(x~;~y~;~z)$.
		
Exprimer MC$^2$ et MD$^2$ en fonction des coordonnées de M.

En déduire qu’une équation cartésienne du plan $P$ est : $-3x - 8y + z = 10$.
		\item Justifier que les plans $P_1$ et $P_2$ sont sécants.
	\end{enumerate}
\item Soit $\Delta$ la droite dont une représentation paramétrique est :
\[\left\{\begin{array}{l c l}
x& =& -2 - 1,9t\\
y&= &\phantom{-2 - 1,9}t\\
z&= & \phantom{-}4 + 2,3t
\end{array}\right.\: \text{où}\: t \in \R\]

Démontrer que $\Delta$ est la droite d’intersection de $P_1$ et $P_2$.
\end{enumerate}

On note $P_3$ le plan médiateur du segment [AC].

On admet qu'une équation cartésienne du plan $P_3$ est : $8x - 10y - 4z= - 15.$

\begin{enumerate}[resume]
\item Démontrer que la droite $\Delta$ et le plan $P_3$ sont sécants.
\item Justifier que le point d’intersection entre $\Delta$ et $P_3$ est le point H.
\end{enumerate}
